\documentclass[12pt,a4paper]{article}
\usepackage[T1]{fontenc}
\usepackage[utf8]{inputenc}
\usepackage[french]{babel}
\usepackage{geometry}
\usepackage{enumitem}
\setlist[itemize]{leftmargin=*,nosep}

\title{Notes de discours --- Présentation orale Mission 4}
\author{Yannis GONOS}
\date{Iveco Group / 3IL Nantes}

\begin{document}
\maketitle

\noindent\textit{Durée totale cible : 9 min 15 s + 45 s de marge = 10 min.}

\section*{Introduction (45 s)}

\begin{itemize}
  \item Accroche : une entreprise industrielle qui emploie des dizaines de milliers de personnes est jugée aujourd'hui autant sur ses KPI de production que sur ses indicateurs extra-financiers.
  \item Objectifs de la mission : comprendre comment la politique RSE d'Iveco Group est structurée, où elle est performante et quelle place le numérique responsable peut y occuper depuis le site de Haute-Goulaine.
  \item Annonce du plan : rapide tour d'horizon de l'entreprise, puis trois parties --- gouvernance et climat, DEI/social, numérique responsable --- avant de conclure sur les perspectives.
\end{itemize}

\section*{Présentation entreprise (1 min)}

\begin{itemize}
  \item Iveco Group : acteur international des véhicules industriels, avec des marques de référence et un fort ancrage européen.
  \item Site de Haute-Goulaine : équipe IT de \textbf{5} personnes en charge du développement logiciel interne.
  \item Environnement : postes Windows 11 / WSL2, stack PHP/Symfony, bases MySQL/Informix, CI/CD et GitLab --- un contexte typique d'une unité agile intégrée au groupe.
  \item On y reviendra en Partie 3, car c'est notre terrain d'action privilégié pour le numérique responsable.
\end{itemize}

\section*{Partie 1 --- Gouvernance et climat (2 min)}

\begin{itemize}
  \item La politique RSE repose sur quatre priorités : réduire l'empreinte carbone, adopter une approche cycle de vie, sécuriser les collaborateurs, et renforcer l'inclusion.
  \item \textbf{Net zéro carbone d'ici 2040} : Iveco s'est engagé dix ans avant l'échéance de Paris. Objectifs intermédiaires : \textbf{-50\,\%} Scope 1 et 2 d'ici 2030 ; \textbf{-38\,\%} d'intensité Scope 3 d'ici 2030.
  \item Le chiffre qui résume tout : \textbf{99,5\,\%} des émissions se situent en Scope 3, c'est-à-dire principalement pendant l'usage des véhicules vendus. Cela montre que le défi carbone du groupe, c'est la phase d'usage.
  \item Le mix énergétique atteint \textbf{35,2\,\%} de renouvelables, en hausse grâce aux achats d'électricité verte.
  \item Points forts ESG : S\&P Global, EcoVadis Gold, CDP Climat A. Point de vigilance : \textbf{CDP Forêts = C}, symptomatique d'une dépendance forte aux matériaux critiques et d'une traçabilité amont encore incomplète.
  \item En gouvernance éthique, la Compliance Helpline a reçu 241 signalements en 2025 et la formation au Code de Conduite couvre environ 12 200 salariés.
\end{itemize}

\section*{Partie 2 --- DEI et social (1 min 45)}

\begin{itemize}
  \item La gouvernance DEI est structurée et s'appuie sur une certification EDGE. Les sujets handicap, diversité culturelle et égalité de rémunération sont pilotés au niveau groupe.
  \item Trois chiffres à retenir :
  \begin{itemize}[label=--]
    \item \textbf{-51\,\%} de taux de blessures entre 2019 et 2025, déjà au-delà de l'objectif \textbf{-40\,\%} fixé pour 2026 ;
    \item \textbf{30\,\%} de femmes dans l'effectif visé d'ici 2028 ;
    \item \textbf{150 000} bénéficiaires touchés par les actions communautaires, encadrées par le référentiel B4SI.
  \end{itemize}
  \item Bilan : des résultats solides, mais des trajectoires DEI à maintenir sur la durée, notamment en diversité de genre et dans l'inclusion du handicap.
\end{itemize}

\section*{Partie 3 --- Numérique responsable (3 min)}

\subsection*{Constat (1 min)}

\begin{itemize}
  \item Point d'entrée : le poste de l'apprenti, avec WSL, CI/CD, revues de code et environnement virtualisé. Des atouts existants.
  \item Constat de terrain personnel : en tant que sous-site rattaché à Guyancourt, Haute-Goulaine ne dispose d'\textbf{aucun relais RSE formalisé} --- pas de référent, pas d'indicateur local. Un site tertiaire de cinq personnes passe mécaniquement sous le radar du déploiement groupe.
  \item Pour autant, des pratiques vertueuses existent déjà de manière informelle : mutualisation par virtualisation, revues de code, faible taux de déplacement. L'enjeu : \textbf{formaliser et mesurer l'existant}, pas créer une politique de toutes pièces.
  \item La politique de confidentialité intègre déjà le \textit{Privacy by Design} ; on peut renforcer cette démarche par des \textbf{PIA} systématisées sur les nouveaux projets IT.
  \item Iveco dispose de briques numériques solides : Web Academy, IMDS, ESGeo.
  \item \textbf{Mais} : à ce jour, aucun document ne formalise un programme \textit{Green IT} appliqué à l'informatique interne : postes, serveurs, cycle de vie du matériel, éco-conception logicielle.
\end{itemize}

\subsection*{Contexte sectoriel (45 s)}

\begin{itemize}
  \item Volvo structure un programme \textit{Green by Design} autour du cloud et du matériel circulaire.
  \item Scania mesure l'impact de ses applications en \textit{vCPU hours} et allonge la durée de vie de son matériel.
  \item Daimler a réduit de \textbf{43\,\%} ses Scope 1 et 2 depuis 2021.
  \item Iveco a les briques, mais ne communique pas encore sur un programme Green IT interne ; c'est justement notre opportunité.
\end{itemize}

\subsection*{Recommandations (1 min 15)}

\begin{itemize}
  \item Recommandation 1 --- \textbf{éco-conception logicielle} : checklist Green IT dans les revues de code, optimisation des requêtes, extinction automatique des environnements non productifs. KPI : \textbf{-15\,\%} de consommation CPU/mémoire des serveurs de développement d'ici fin 2027. Argument fort : GitLab et les pipelines CI/CD sont hébergés \textbf{sur site}, donc la mesure est immédiatement réalisable avec une supervision légère (Prometheus/Grafana), sans coût de licence ni dépendance à un prestataire cloud.
  \item Recommandation 2 --- \textbf{cycle de vie du matériel} : privilégier le réemploi/reconditionnement et allonger la durée d'usage. KPI : \textbf{30\,\%} de parc test/reconditionné d'ici 2028.
  \item Recommandation 3 --- \textbf{achats responsables} : introduire EPEAT, TCO Certified et rapports de durabilité fournisseurs. KPI : \textbf{100\,\%} des nouveaux postes labellisés d'ici 2026.
  \item Ces objectifs sont réalistes pour une équipe de cinq personnes, à condition d'établir une base de référence au premier trimestre 2026 et de prioriser deux ou trois leviers sans se disperser.
\end{itemize}

\section*{Conclusion (45 s)}

\begin{itemize}
  \item Gouvernance RSE mûre, avec un défi structurant : le Scope 3, c'est-à-dire l'usage des véhicules vendus.
  \item Résultats sociaux solides, mais trajectoires DEI à maintenir.
  \item Le numérique responsable est le \textbf{levier d'action le plus personnel} : mesurer, prioriser et piloter avec 2--3 KPI concrets.
  \item Message final : au-delà du constat, une feuille de route opérationnelle pour l'équipe IT de Haute-Goulaine, alignée sur la stratégie RSE du groupe.
\end{itemize}

\vfill
\noindent\textit{Conseil d'animation : parler sans lire, faire le lien entre le chiffre affiché et sa signification, et anticiper les 45 secondes de marge pour les questions.}

\end{document}
